% Pollis New Methods: random forests (out-of-bag error and permutation importance) — draft 1 (11 October 2026)
% Compile with: latexmk -pdf nm_random_forests.tex && latexmk -c
% Shared look for the Pollis New Methods documents (same fonts and colours as the brochure and the workflow handouts).
\documentclass[10pt,a4paper]{article}

\usepackage[a4paper,left=18mm,right=18mm,top=15mm,bottom=17mm,footskip=8mm]{geometry}
\usepackage[T1]{fontenc}
\usepackage[default,semibold]{sourcesanspro}
\usepackage{textcomp}
\usepackage[italic]{mathastext}
\usepackage{amsmath}
\usepackage{microtype}
\usepackage[table]{xcolor}
\usepackage{tikz}
\usepackage[most]{tcolorbox}
\usepackage{fontawesome5}
\usepackage{tabularx,booktabs}
\usepackage{enumitem}
\usepackage{fancyhdr}
\usepackage[hidelinks]{hyperref}

% ---------------------------------------------------------------- colours (as in the brochure)
\definecolor{navy}{HTML}{1C3A5E}
\definecolor{pblue}{HTML}{3B78B5}
\definecolor{amber}{HTML}{F2B33D}
\definecolor{deepamber}{HTML}{D9902A}
\definecolor{pale}{HTML}{EEF3F9}
\definecolor{ink}{HTML}{2A2F36}
\definecolor{logonavy}{HTML}{173D63}
\definecolor{logodark}{HTML}{C98A1E}
\definecolor{logoamber}{HTML}{F2B446}
\definecolor{logocream}{HTML}{FEEFD2}
\definecolor{soft}{HTML}{5E6773}
\definecolor{agree}{HTML}{2F7D4F}
\definecolor{differ}{HTML}{B23B3B}
\color{ink}

\setlength{\parindent}{0pt}
\setlength{\parskip}{4pt}
\setlist{nosep,leftmargin=*,itemsep=2pt}
\renewcommand{\arraystretch}{1.15}

\pagestyle{fancy}
\fancyhf{}
\renewcommand{\headrulewidth}{0pt}
\fancyfoot[L]{\footnotesize\color{soft}\textbf{Pollis}\enspace New Methods\enspace·\enspace\nmtheme}
\fancyfoot[R]{\footnotesize\color{soft}\thepage}
\newcommand{\nmtheme}{}
\newcommand{\nmdate}{11 October 2026}

% The Pollis logo, redrawn so it prints sharply at any size
\newcommand{\pollislogo}[1]{%
  \begin{tikzpicture}[x=#1/512,y=-#1/512,baseline=(current bounding box.center)]
    \fill[logonavy] (0,0) rectangle (512,512);
    \fill[logodark] (36,422) -- (104,256) -- (172,422) -- cycle;
    \fill[logoamber] (104,422) -- (203,180) -- (302,422) -- cycle;
    \fill[logocream] (203,422) -- (339,89) -- (476,422) -- cycle;
  \end{tikzpicture}}

% Top of the first page: logo, name; then the kicker, title and a one-line promise
\newcommand{\nmhead}[3]{%
  \renewcommand{\nmtheme}{#1}%
  \thispagestyle{fancy}%
  \begin{minipage}[c]{13mm}\pollislogo{11mm}\end{minipage}%
  \begin{minipage}[c]{\dimexpr\linewidth-13mm\relax}
    {\color{navy}\Large\bfseries Pollis}\enspace{\color{soft}New Methods}\hfill
    {\color{deepamber}\small\bfseries\MakeUppercase{#1}}
  \end{minipage}\par\vspace{7mm}
  {\color{navy}\fontsize{22}{26}\selectfont\bfseries #2}\par\vspace{4pt}
  {\color{soft}\large #3}\par\vspace{5mm}}

\newcommand{\nmsection}[1]{\par\vspace{5pt}{\color{navy}\large\bfseries #1}\par\vspace{1pt}}
\newcommand{\nmsub}[1]{\par\vspace{2pt}{\color{deepamber}\bfseries #1}\par}
\newcommand{\gt}{\,{\color{deepamber}\bfseries ›}\,}
\newcommand{\menupath}[1]{{\bfseries\color{navy}#1}}
\newcommand{\fn}[1]{\texttt{\small #1}}
% Citations: \nmcite{key}{Breiman (2001)} links to the entry \nmref{key}{...} in the References
\newcommand{\nmcite}[2]{\hyperlink{ref:#1}{\color{pblue}#2}}
\newcommand{\nmref}[2]{\raisebox{1.1em}[0pt][0pt]{\hypertarget{ref:#1}{}}#2\par}
\newcommand{\nmdoi}[1]{\href{https://doi.org/#1}{\color{pblue}doi:#1}}
% Where a notebook is in Pollis: its title in Browse Methods, and its file
\newcommand{\browse}[2]{\menupath{Compose \gt New Methods \gt Browse Methods}: \emph{#1} (\fn{#2.ipynb})}

% Key facts at the top: source, missing in Julia, notebook, result
\newcommand{\factsheet}[4]{%
  \begin{tcolorbox}[enhanced, colback=pale, colframe=pale, boxrule=0pt, sharp corners,
    left=7pt, right=7pt, top=5pt, bottom=5pt, fontupper=\small]
  \begin{tabularx}{\linewidth}{@{}>{\leavevmode\bfseries\color{navy}}p{27mm}>{\leavevmode}X@{}}
  Tested against & #1\\
  Missing in Julia & #2\\
  Notebook & #3\\
  Result & #4
  \end{tabularx}
  \end{tcolorbox}}

% Boxes
\newtcolorbox{pbox}[2][]{enhanced, colback=pale, colframe=amber, boxrule=0pt,
  leftrule=3pt, sharp corners, left=7pt, right=7pt, top=6pt, bottom=6pt,
  fontupper=\small\raggedright, before upper={{\color{navy}\normalsize\bfseries #2}\par\smallskip}, #1}
\newtcolorbox{nbox}[2][]{enhanced, colback=navy, colframe=navy, boxrule=0pt,
  sharp corners, left=8pt, right=8pt, top=7pt, bottom=7pt, coltext=white,
  fontupper=\small\raggedright, before upper={{\color{amber}\normalsize\bfseries #2}\par\smallskip}, #1}

% Value tables: published beside Julia
\newenvironment{valuetable}[1]{%
  \par\small\tabularx{\linewidth}{#1}\toprule}{\bottomrule\endtabularx\par}
\newcommand{\ok}{{\color{agree}\faCheck}}
\newcommand{\no}{{\color{differ}\faTimes}}

% End of the document: credit and licence
\newcommand{\nmcredit}{%
  \par\vfill
  \begin{minipage}[b]{0.78\linewidth}\footnotesize\color{soft}
  Methods written and tested with Claude Code (Anthropic) in Pollis, from the published papers and their data.\par
  \textcopyright\ 2026 Antonio Saragga Seabra\enspace·\enspace Licensed under \href{https://creativecommons.org/licenses/by-sa/4.0/}{CC BY-SA 4.0}\enspace·\enspace Draft, \nmdate
  \end{minipage}\hfill
  \begin{minipage}[b]{0.18\linewidth}\raggedleft\pollislogo{9mm}\par\vspace{2pt}{\footnotesize contact@pollis.org}\end{minipage}}


\begin{document}

\nmhead{Random Forests}
  {Out-of-bag error and variable importance for random forests}
  {Breiman's own measures of a forest, tested against his Table 2 on 19 data sets.}

\factsheet
  {\nmcite{breiman2001}{Breiman (2001)}, \emph{Random forests}, \emph{Machine Learning}: Table 2 (p.\,11) and Section 10 (pp.\,19--21).
   14 real data sets from PollisDatasets, four simulated as Breiman did, and the zip-code digits of \nmcite{hastie2009}{Hastie, Tibshirani and Friedman (2009)}
   (downloaded only if the user agrees).}
  {the out-of-bag error and permutation importance. DecisionTree.jl grows forests, but has neither.}
  {\browse{Out-of-bag error and variable importance for random forests}{oob}}
  {34 of 38 Selection and Single input errors within simulation noise of Table 2; importance picks the inputs Breiman's did.}

\nmsection{The problem}
A random forest grows each tree on a bootstrap sample, so about a third of the cases are left out of each tree. \nmcite{breiman2001}{Breiman (2001)} used
these out-of-bag cases twice: to estimate the forest's error without a test set, and to measure how much the forest relies on
each input. Both are cited constantly and computed differently by different libraries. Breiman's values are averages of 100
random runs, so the test is statistical: a Julia average, with its standard error, against his.

\nmsection{The method}
\nmsub{Forest-RI and the out-of-bag error}
Forest-RI (p.\,9) grows each tree on a bootstrap sample of the training cases, splitting at each node on the best of $F$
inputs drawn at random, by the Gini criterion of \nmcite{breiman1984}{Breiman, Friedman, Olshen and Stone (1984)}, unpruned; the forest votes. Run down only the trees that did not see them, the
out-of-bag cases give the forest's \textbf{out-of-bag error} (Section 3.1), and the error of each tree alone, Table 2's
``One tree''. \fn{forest} returns the trees and both errors; \fn{classify} runs cases down a tree without recursion, since the
deepest trees overflow the stack of DecisionTree.jl's recursive \fn{apply\_tree}.
\nmsub{Permutation importance}
Permute input $m$ among the out-of-bag cases of each tree and run them down again: the percent increase in the out-of-bag
error measures how much the forest relies on $m$ (Section 10).

\nmsection{What the tests show}
Breiman grew two forests per data set, with $F = 1$ and $F = \mathrm{int}(\log_2 M + 1)$, and took the test error of the one
with the lower \emph{out-of-bag} error: the \textbf{Selection} column. \textbf{Single input} is the test error with $F = 1$.
A value is \emph{within noise} when the two averages differ by less than two standard errors of a difference ($\sqrt2$ s.e.).
\begin{valuetable}{@{}lXl@{}}
\textbf{Test} & \textbf{What is compared} & \textbf{Within noise}\\\midrule
Table 2, Selection & test error of the forest chosen by out-of-bag error & 17 of 19\\
Table 2, Single input & test error of the forest with $F = 1$ & 17 of 19\\
Table 2, One tree & out-of-bag error of single trees & 7 of 19\\
Section 10 & errors of forests on the inputs importance picks & 4 of 5\\
\end{valuetable}

\nmsub{Where they differ}
Liver Selection 27.3 vs 25.1 (2.0 s.e.); twonorm Single input 3.7 vs 3.9 (2.1 s.e.); ringnorm Selection 4.4 vs 4.9 and Single
input 4.3 vs 4.9 (2.4 and 4.7 s.e.). Both ringnorm values are lower than Breiman's.

One tree: most differ by about one point (pima 31.9 vs 33.1, liver 39.3 vs 40.6), and vowel by more, 36.1 vs 30.4. The test is
strict there: each value averages 100 trees over 100 runs, so the standard errors are small, and small differences between
DecisionTree.jl's tree code and Breiman's (tie-breaking, split points) show as several standard errors.

\nmsub{Importance picks the same inputs}
On diabetes, importance ranks input 2 (plasma glucose) first, as Breiman's Figure 4 does. Forests on the inputs it picks give
29.9 / 24.0 / 30.2 (input 2 / all eight / 2 and 8) against 29.7 / 23.1 / 29.4, within noise, and 28.1 vs 26.4 for inputs 2
and 6 (2.4 s.e.). On the 1984 House votes, input 4 (freezing physicians' fees) comes first, as in Figure 6, and a forest on it
alone errs 4.7\% against 4.3\%, within noise.

\nmsection{Still open}
\begin{itemize}
  \item \textbf{Ringnorm and twonorm.} The forests are below Breiman on these two simulated sets, and arc-fs (its own document)
    shows the same small gap on the same two sets. A gap shared by two methods points to how the trees are grown, or to the
    generators, which follow the description in \nmcite{breiman1998}{Breiman (1998)}, p.\,7. Not explained.
  \item \textbf{One tree}, about one point off on most data sets: a difference in the single trees, not in the forests' votes.
\end{itemize}

\nmsection{Further considerations}
\begin{itemize}
  \item \textbf{Replicating random results.} Every run has its own seed, so the results do not depend on the number of
    threads, and every comparison carries its standard error: a disagreement is a number of standard errors, not a judgement.
  \item \textbf{Upstream.} Out-of-bag error and permutation importance are small additions to DecisionTree.jl, and Table 2
    is a ready-made benchmark for them.
  \item \textbf{Next.} Breiman's third use of out-of-bag cases, proximities (for outliers and missing values), and his
    regression forests (Section 11), each tested against his tables.
\end{itemize}

\nmsection{References}
{\footnotesize
\nmref{breiman1998}{Breiman, L. (1998). Arcing classifiers. \emph{The Annals of Statistics} 26(3), 801--849. \nmdoi{10.1214/aos/1024691079}}
\nmref{breiman2001}{Breiman, L. (2001). Random forests. \emph{Machine Learning} 45(1), 5--32. \nmdoi{10.1023/A:1010933404324}}
\nmref{breiman1984}{Breiman, L., Friedman, J.\,H., Olshen, R.\,A. and Stone, C.\,J. (1984). \emph{Classification and Regression Trees}. Wadsworth.}
\nmref{hastie2009}{Hastie, T., Tibshirani, R. and Friedman, J. (2009). \emph{The Elements of Statistical Learning}, 2nd ed. Springer: the zip-code data.}
}

\nmcredit
\end{document}
